\subsection{一般线性椭圆问题}\label{sec:general-linear-elliptic-problems}

本章至今只讨论了椭圆问题的 Dirichlet 边界条件. 本节讨论物理与数学问题中出现的其他几种边界条件. 由于物理问题提出了种类繁多的边界条件, 这里的讨论并不穷尽所有情形.

\subsubsection{Neumann 问题}\label{subsec:neumann-problem}

在 Dirichlet 问题之后, 第二种最常见且重要的椭圆边值问题是 Neumann 问题.

\begin{Definition}
\label{def:classical-neumann-solution}
设 $\Omega\subset\mathbb{R}^n$ 为具有 $C^1$ 边界的有界区域, 给定 $f\in C_b(\Omega)$. 若函数\[u\in C_b^2(\Omega)\cap C_b^1(\overline\Omega)\]在 $\Omega$ 中满足
\begin{equation}\label{eq:neumann-classical-divergence-equation}
L(\bm{x},D)u:=\sum_{i,j=1}^n\frac{\partial}{\partial x_i}a_{ij}(\bm{x})\frac{\partial u}{\partial x_j}+\sum_{i=1}^nb_i(\bm{x})\frac{\partial u}{\partial x_i}+c(\bm{x})u=f
\end{equation}
且在 $\partial\Omega$ 上满足
\begin{equation}\label{eq:neumann-homogeneous-conormal-boundary-condition}
\sum_{i,j=1}^na_{ij}(\bm{x})\frac{\partial u}{\partial x_j}\eta_i(\bm{x})=0
\end{equation}
则称其为 Neumann 问题的\textbf{经典解}, 其中 $\bm{\eta}(\bm{x})$ 为 $\partial\Omega$ 在 $\bm{x}$ 处的单位外法向量.
\end{Definition}

与 Dirichlet 问题一样, 可以定义 Neumann 问题的强解.

\begin{Definition}
\label{def:strong-neumann-solution}
设 $\Omega\subset\mathbb{R}^n$ 为具有 $C^1$ 边界的有界区域, 给定 $f\in L^2(\Omega)$. 若函数\[u\in H^2(\Omega)\]在 $L^2(\Omega)$ 中满足 \eqref{eq:neumann-classical-divergence-equation}, 且在 $\partial\Omega$ 上按迹的意义满足 \eqref{eq:neumann-homogeneous-conormal-boundary-condition}, 则称其为 Neumann 问题的\textbf{强解}.
\end{Definition}

为将 Neumann 问题陈述为弱形式, 如前所述利用分部积分, 从微分算子 $L$ 和边界条件构造双线性型. 注意, 对 $H^2(\Omega)$ 中的任意 $\phi$ 和 $u$, 有
\begin{equation}\label{eq:neumann-integration-by-parts-boundary-form}
\begin{aligned}
\int_\Omega\phi Lu\,\mathrm{d}\bm{x}&=B[\phi,u]\\
&\quad+\sum_{i,j=1}^n\int_{\partial\Omega}a_{ij}(\bm{x})\frac{\partial u}{\partial x_j}(\bm{x})\eta_i(\bm{x})\phi(\bm{x})\,\mathrm{d}\bm{x},
\end{aligned}
\end{equation}
其中双线性型 $B$ 在 \eqref{eq:second-order-garding-bilinear-form} 中定义. 因而, 若 $u$ 满足边界条件 \eqref{eq:neumann-homogeneous-conormal-boundary-condition}, 则对每个 $\phi\in H^1(\Omega)$ 有
\begin{equation}\label{eq:neumann-natural-variational-integration-identity}
\int_\Omega\phi Lu\,\mathrm{d}\bm{x}=B[\phi,u]
\end{equation}
事实上, 将此作为 Neumann 问题弱解的定义.

\begin{Definition}
\label{def:weak-neumann-solution}
设 $\Omega\subset\mathbb{R}^n$ 为有界区域, 给定 $f\in L^2(\Omega)$. 若函数\[u\in H^1(\Omega)\]对每个 $v\in H^1(\Omega)$ 均满足
\begin{equation}\label{eq:weak-neumann-variational-equation}
B[v,u]=(f,v)
\end{equation}
则称其为 Neumann 问题的\textbf{弱解}.
\end{Definition}

有必要作几点说明.\begin{enumerate}\item 当然, 上述构造保证 Neumann 问题的每个强解也都是弱解.\item 在弱形式的构造中, 边界条件 ``消失\textquotesingle\textquotesingle\ 了; 即条件 \eqref{eq:neumann-homogeneous-conormal-boundary-condition} 既不显式出现在双线性型中, 也不显式出现在容许函数空间中. 因而称 Neumann 条件为\emph{自然边界条件}. 事实上, 由于迹定理不保证 $H^1(\Omega)$ 函数的法向导数存在, 在 \eqref{eq:weak-neumann-variational-equation} 的弱解上计算边界条件 \eqref{eq:neumann-homogeneous-conormal-boundary-condition} 未必有意义.\item 假设了数据 $f$ 属于空间 $L^2(\Omega)$. 这一条件可以减弱, 例如可取 $f\in L^2(S)$, 其中 $S$ 为包含在 $\Omega$ 中的光滑曲面. 然而, 不能像 Dirichlet 问题那样取 $H^{-1}(\Omega)$ 中的任意数据.\end{enumerate}
